\section{Future Plans}

There are a few things I would do differently or extend if I were to continue working
on this project.

The most obvious next step is completing the GitOps loop that ArgoCD was supposed to
close. The infrastructure repository and the GitHub Actions workflow are already in place.
The cleanest path forward given the Sandbox limitations is to replace the ArgoCD sync
step with a direct \texttt{oc apply} in the Actions workflow, using a service account
token scoped to the project namespace. This would give the same end result, a push to
the application repo triggers a deployment to the cluster automatically, without needing
cluster-admin access.

For the secrets management I used hardcoded placeholder values throughout this project,
which is acceptable for a lab environment but would not be acceptable in production. A
proper solution would use something like HashiCorp Vault or the OpenShift Secrets Store
CSI Driver, so that credentials are generated, rotated, and injected without any human
ever handling the actual values.

I would also add health checks and readiness probes to the backend Deployment. I added
one to the frontend to fix the Route issue, but the backend is currently relying on the
cluster just waiting long enough for it to be ready. An explicit readiness probe pointing
at the \texttt{/health} endpoint would make the rollout more reliable, especially when
the backend is waiting for PostgreSQL to become available.

On the monitoring side, the application has no observability beyond the basic container
logs. Adding a metrics endpoint and connecting it to the built-in OpenShift monitoring
stack would make it much easier to understand how the application is behaving in
production.
